As shown in  \cref{thm:gs-mdss-correctness} and \cref{thm:gs-mdss-ul}, a suitable choice of field is required to achieve unlinkability and MD-correctness.
Indeed, if a dealer randomly samples the same evaluation point twice and produces two identical shares, then unlinkability is compromised.
While this will occur with negligible probability if fields are large enough, in practice small fields are often best to minimize bandwidth and decoding complexity.
We therefore propose a stateful variant of our construction in which dealers keep track of previously sampled evaluation points and emit {\em noise shares} (comprising random elements in place of the polynomial evaluation) whenever an $x$-coordinate is repeated.
This preserves unlinkability at the cost of decreasing the number of ``useful'' shares passed to the reconstruction algorithm.
This in turn results in a non-negligible probability of failure, which may be an acceptable tradeoff for protocols in which guaranteed recovery is not necessary.
A formal description of this scheme can be seen in \cref{fig:gs-mdss-small-field}.

\begin{figure}
	\begin{pcvstack}[boxed, center, space=1em]
		\procedure[linenumbering]{$\mdssShareAlg(\secparam, \mdssSecret, \mdssNumShares)$}{
			\text{Sample a polynomial } p \text{ of degree } \mdssPrivThr \text{ where } \mdssSecret = p(0) \\
			\text{Sample } \mdssNumShares \text{ field elements } x_1, \dots, x_\mdssNumShares \sample \FF \\
			\pcfor i \in [\mdssNumShares] \pcdo \\
			\t \pcif x_i \notin \set{x_1,\dots,x_{i - 1}} \pcthen \\
			\t \t \mdssBlindShare_i := (x_i, p(x_i)) \\
			\t \pcelse \\
			\t \t y_i \sample \FF \\
			\t \t \mdssBlindShare_i := (x_i, y_i) \\
			\pcreturn \set{\mdssBlindShare_i}_{i \in [\mdssNumShares]}
		}

		\procedure[linenumbering]{$\mdssReconstructAlg(\mdssBlindShare_1, \dots, \mdssBlindShare_{\mdssWindowSize})$}{
			p_1,\dots,p_{\mdssNumShareSets} \gets \gsDecode(\mdssPrivThr, \mdssRecThr, \set{\mdssBlindShare_1, \dots, \mdssBlindShare_{\mdssWindowSize}}) \\
			\pcreturn \set{p_1(0), \dots, p_{\mdssNumShareSets}(0)}
		}
	\end{pcvstack}
	\caption{An MDSS construction for small fields $\sfmdss$}
	\label{fig:gs-mdss-small-field}
\end{figure}

\newcommand{\challengeShare}{\ensuremath{\mdssShare{\mdssSecret_b}^*}}

\begin{theorem}
	For all finite fields $\FF$, the construction shown in \cref{fig:gs-mdss-small-field} satisfies $\mdssULVar$-unlinkability for $\mdssULVar(\secpar) = 0$.
\end{theorem}
\begin{proof}
	We begin by showing that $\sfmdss$ satisfies one-more unlinkability.
	Consider the distribution of $\challengeShare$ in the one-more unlinkability game.
	Note that whether $b = 0$ or $b = 1$ the $x$-coordinate of $\challengeShare$ is a uniformly random element of $\FF$.
	We therefore limit our analysis to its $y$-coordinate.
	If $b = 0$ and the $x$-coordinate of $\challengeShare$ collides with a previous share, then its $y$-coordinate is a uniformly random element of $\FF$.
	If $b = 0$ and the $x$-coordinate of $\challengeShare$ does \emph{not} collide, then it is at most the $\mdssPrivThr$'th point on a degree $\mdssPrivThr +1$ polynomial, and so its $y$-coordinate is also a uniformly random element of $\FF$.
	Finally, when $b = 1$, we again have that the $y$-coordinate is a uniformly random element of $\FF$.
	Therefore, the distribution in both games is identical, and the one-more unlinkability advantage of all adversaries must be $0$.

	The theorem then follows directly from \cref{thm:omul-implies-ul}.
\end{proof}

\begin{theorem}
	For all sets of parameters satisfying ..., and finite fields $\FF$ such that $\mdssMaxShares \in O(\secpar^d)$ and $|\FF| \in O(\secpar^{d'})$ for constants $d$, $d'$, the construction show in \cref{fig:gs-mdss-small-field} satisfies $\mdssCVar$-correctness for $\mdssCVar(\secpar) = \dots$.
\end{theorem}

